\subsection{MONOID OPERATING ON A SET}

DEFINITION 1. Let M be a monoid, with law written multiplicatively and identity element denoted by e, and E a set. An action \(\alpha \mapsto f_{\alpha}\) of M on E is called a left (resp. right) operation of M on E if \(f_{e} = \mathrm{Id}_{\mathrm{E}}\) and \(f_{\alpha \beta} = f_{\alpha} \circ f_{\beta}\) (resp. \(f_{\alpha \beta} = f_{\beta} \circ f_{\alpha}\) ) for all \(\alpha, \beta \in \mathbf{M}\) .

In other words, a left (resp. right) operation of a monoid M on a set E is a monoid homomorphism of M into the monoid \(E^{E}\) (resp. the opposite monoid of \(E^{E}\) ) with composition of mappings. If the law of action corresponding to the action of M is denoted by left (resp. right) multiplication, the fact that this action is a left (resp. right) operation may be expressed by the formulae

\[
e. x = x; \alpha . (\beta . x) = (\alpha \beta). x \quad \text { for } \alpha , \beta \in \mathrm{M} \text { and } x \in \mathrm{E}.\tag{1}
\]

\[
(\text { resp. } x. e = x; (x. \alpha). \beta = x. (\alpha \beta) \quad \text { for } \alpha , \beta \in M \text { and } x \in E).
\]

Under these conditions, it is also said that M operates on E on the left (resp. right) and that the corresponding laws of action are laws of left (resp. right) operation of the monoid M on E.

Let M be a monoid; a set E with a left (resp. right) operation of M on E is called a left (resp. right) M-set. The monoid M is said to operate on the left (resp. right) faithfully if the mapping \(\alpha \mapsto f_{\alpha}\) of M into \(E^{E}\) is injective.

Examples. (1) Let E be a set; the canonical action of \(E^{E}\) on E (§ 3, no. 1, Example 3) is a left operation.

\begin{enumerate}
\def\labelenumi{(\arabic{enumi})}
\setcounter{enumi}{1}
\tightlist
\item
  Let M be a monoid. The left (resp. right) action of M on itself derived from the law on M (§ 3, no. 3, Example 7) is a left (resp. right) operation of M on itself. When considering this operation, we say that M operates on itself by left (resp. right) translation.
\end{enumerate}

Let E be a left (resp. right) M-set and \(M^{0}\) the opposite monoid to M. Under the same action, the monoid \(M^{0}\) operates on E on the right (resp. left). The \(M^{0}\) -set obtained is called opposite to the M-set E. The definitions and results relating to left M-sets carry over to right \(M^{0}\) -sets when passing to the opposite structures.

In the rest of this paragraph, we shall consider, unless otherwise stated, only left M-sets which we shall call simply M-sets. Their law of action will be denoted by left multiplication.

Let E be a set. Let G be a group operating on E. For all \(\alpha\) in G, the element of \(E^{E}\) defined by \(\alpha\) is a permutation of E ( \(\S\) 2, no. 3, Example 2). Being given an operation of G on E therefore amounts to being given a homomorphism of G into \(S_{E}\) .

In conformity with § 3, no. 3, we make the following definition:

DEFINITION 2. Let M be a monoid and E and E' M-sets. A mapping f of E into E' such that, for all \(x \in E\) and all \(\alpha \in M\) , \(f(\alpha \cdot x) = \alpha \cdot f(x)\) is called an M-set homomorphism (or M-morphism, or mapping compatible with the operations of M).

The identity mapping of an M-set is an M-morphism. The composition of two M-morphisms is an M-morphism. For a mapping of one M-set into another to be an isomorphism, it is necessary and sufficient that it be a bijective M-morphism and the inverse mapping is then an M-morphism.

Let \((\mathbf{E}_i)_{i\in \mathbf{I}}\) be a family of M-sets and let \(\mathbf{E}\) be the product set of \(\mathbf{E}_i\) . The monoid M operates on \(\mathbf{E}\) by \(\alpha .(x_i)_{i\in \mathbf{I}} = (\alpha .x_i)_{i\in \mathbf{I}}\) and \(\mathbf{E}\) , with this action, is an M-set; let \(\mathbf{E}'\) be an M-set; a mapping \(f\) of \(\mathbf{E}'\) into \(\mathbf{E}\) is an M-morphism if and only if \(\mathbf{pr}_i\circ f\) is an M-morphism of \(\mathbf{E}'\) into \(\mathbf{E}_i\) for all \(i\in \mathbf{I}\) .

Let E be an M-set and F a stable subset of E under the action of M; with the induced law, F is an M-set and the canonical injection \(F \rightarrow E\) is an M-morphism.

Let E be an M-set and R an equivalence relation on E compatible with the action of M; the quotient E/R with the quotient action is an M-set and the canonical mapping \(E \rightarrow E/R\) is an M-morphism.

Let \(\phi: \mathbf{M} \to \mathbf{M}'\) be a monoid homomorphism, \(\mathbf{E}\) an \(\mathbf{M}\) -set and \(\mathbf{E}'\) an \(\mathbf{M}'\) -set. A mapping \(f\) of \(\mathbf{E}\) into \(\mathbf{E}'\) such that, for all \(x \in \mathbf{E}\) and \(\alpha \in \mathbf{M}\) ,

\[
f (\alpha . x) = \phi (\alpha). f (x)
\]

is called a \(\phi\) -morphism of E into E' (cf.~§ 3, no. 1).

Extension of a law of operation. Given (for example) three sets \(F_{1}, F_{2}, F_{3}\) , permutations \(f_{1}, f_{2}, f_{3}\) of \(F_{1}, F_{2}, F_{3}\) respectively and an echelon F on the base sets \(F_{1}, F_{2}, F_{3}\) (Set Theory, IV, § 1, no. 1), we can define, proceeding step by step on the construction of the echelon F, a permutation of F called the canonical extension of \(f_{1}, f_{2}, f_{3}\) to F (Set Theory, IV, § 1, no. 2); we shall denote it by \(\phi_{F}(f_{1}, f_{2}, f_{3})\) .

Then let G be a group and \(h_{i}\) a homomorphism of G into the symmetric group of \(F_{i}\) (i = 1, 2, 3), in other words an operation of G on \(F_{i}\) . The mapping \(x \mapsto x_{F} = \phi_{F}(h_{1}(x), h_{2}(x), h_{3}(x))\) is a homomorphism of G into \(S_{F}\) , in other words an operation of G on F, called the extension of \(h_{1}, h_{2}, h_{3}\) to F. Let P be a subset of F such that, for all \(x \in G\) , \(x_{F}(P) = P\) ; let \(x_{P}\) be the restriction of \(x_{F}\) to

P; then the mapping \(x \mapsto x_{P}\) is an operation of G on P, also called the extension of \(h_{1}, h_{2}, h_{3}\) to P.

For example, let \(\mathbf{K}\) and \(\mathbf{L}\) be two echelons on \(\mathbf{F}_1, \mathbf{F}_2, \mathbf{F}_3\) ; take \(\mathbf{F}\) to be the set of subsets of \(\mathbf{K} \times \mathbf{L}\) and \(\mathbf{P}\) to be the set of mappings of \(\mathbf{K}\) into \(\mathbf{L}\) , identified with their graphs. If \(w \in \mathbf{P}\) and \(x \in \mathbf{G}\) , \(x_{\mathbf{P}}(w)\) is the mapping \(k \mapsto x_{\mathbf{L}}(w(x_{\mathbf{K}}^{-1}(k)))\) of \(\mathbf{K}\) into \(\mathbf{L}\) .

